% Generated draft: requires mathematical and source review.
\documentclass[11pt]{article}
\usepackage[margin=25mm]{geometry}
\usepackage{fontspec}
\setmainfont{Latin Modern Roman}
\usepackage{amsmath,amssymb,mathtools,unicode-math}
\setmathfont{Latin Modern Math}
\usepackage{enumitem,needspace,xurl,hyperref}
\setlength{\parindent}{0pt}
\setlength{\parskip}{5pt}
\setlength{\emergencystretch}{3em}
\newcommand{\N}{\mathbb{N}}
\newcommand{\Z}{\mathbb{Z}}
\newcommand{\Q}{\mathbb{Q}}
\newcommand{\R}{\mathbb{R}}
\newcommand{\C}{\mathbb{C}}
\newcommand{\F}{\mathbb{F}}
\newcommand{\A}{\mathbb{A}}
\newcommand{\PP}{\mathbb{P}}
\newcommand{\E}{\mathbb{E}}
\newcommand{\eps}{\varepsilon}
\newcommand{\dd}{\,\mathrm{d}}
\begin{document}
\section*{UnsolvedMath ID: 30005851}
\textbf{Source problem code:} catalogue entry
\subsection{Definitions and mathematical statement}
The supplied record asks for spinorial approaches to two stability problems and an extension of a third. For each integer \(n\geq 3\), the first questions are whether spinors can establish Geroch stability and stability of the positive mass theorem in dimension \(n\). The further question is whether the stability result designated Theorem 7 in the source report can be extended to the scalar-curvature rigidity settings of Goette and Semmelmann, cited there as [23]. The supplied reconstructed statement identifies this as an extension of Llarull stability.

These are research questions rather than fully specified theorem statements. Here a stability assertion would require a specified class of manifolds, a sequence approaching equality in a specified scalar-curvature rigidity or mass inequality, and a specified conclusion of convergence to an appropriate rigid model. The record does not supply those classes, inequalities, convergence notions, or models for the requested assertions. It also does not specify spin hypotheses, asymptotic conditions for the positive-mass question, boundary conditions for the requested generalizations, or the statements of Theorem 7 and [23]. Consequently, no unique fully quantified stability conjecture can be extracted without additional source information.

The adjoining source passage supplies the following precise background result, called the Volume Above Distance Below theorem and attributed to Allen, Perales, and Sormani [11]. Let \((M^m,g_0)\) be a fixed compact oriented Riemannian manifold without boundary, and let \((g_j)_{j\in\N}\) be continuous Riemannian metrics on the same manifold \(M\). Suppose there are real numbers \(C(j)\downarrow0\) and a finite constant \(D_0\) such that, for every \(j\in\N\), every \(p\in M\), and every \(v\in T_pM\),
\[
 g_j(v,v)\geq (1-C(j))g_0(v,v),
 \qquad \operatorname{Diam}(M,g_j)\leq D_0,
\]
and suppose
\[
 \operatorname{Vol}(M,g_j)\longrightarrow\operatorname{Vol}(M,g_0).
\]
The supplied passage states that these hypotheses imply volume-preserving Sormani--Wenger intrinsic flat convergence,
\[
 d_{\mathrm{VF}}\bigl((M,g_j),(M,g_0)\bigr)\longrightarrow0.
\]
Here \(\operatorname{Diam}\) and \(\operatorname{Vol}\) denote Riemannian diameter and total volume, and \(d_{\mathrm{VF}}\to0\) is the source's notation for intrinsic flat convergence together with convergence of total volumes; its particular distance convention is not supplied. This theorem is quoted background, not a newly asserted solution of the requested spinorial problems.

The passage also attributes a version for manifolds with boundary to Allen and Perales [10], without supplying its hypotheses. Its additional Question 6 asks for a version of this background theorem for sequences whose manifolds are not diffeomorphic to the limit manifold. Such an extension requires a replacement for the displayed comparison of metrics on a common tangent bundle; no replacement is specified. The following sentence, concerning extensions of Theorems 6 and 7 to \(1\)-Lipschitz objects, is truncated in the supplied record, so its intended conclusion cannot be stated precisely.
\subsection{Short English statement}
The main questions ask whether spinors establish Geroch and positive-mass stability in every dimension \(n\geq3\), and whether Llarull stability extends to the Goette--Semmelmann rigidity settings. Essential hypotheses and the meaning of stability are missing. The accompanying metric-comparison theorem is recorded as background; the additional question about changing manifold topology is also incomplete.
\subsection{Sources}
\begin{itemize}
\item Local UnsolvedMath registry: supplied record with permanent ID \(30005851\), title \emph{Spinorial Scalar-Curvature Stability}, exact source problem code \texttt{catalogue entry}, in the Oberwolfach Reports Open Problems collection. No individual proposer is supplied.
\item \emph{Analysis, Geometry and Topology of Positive Scalar Curvature Metrics} (2024), Oberwolfach Reports, pp.\ 483--566, DOI \texttt{10.4171/owr/2024/9}: \url{https://doi.org/10.4171/owr/2024/9}. The supplied excerpt includes the page marker 493. It attributes the background Theorem 8 to Allen, Perales, and Sormani [11], a boundary version to Allen and Perales [10], and the proposed rigidity generalization to results of Goette and Semmelmann [23]. Full bibliographic entries for [10], [11], and [23] and the text of Theorem 7 are not supplied.
\item Supplied literature source: \emph{Stability of the positive mass theorem under Ricci bounds}, \url{https://arxiv.org/abs/2111.05202}.
\item Supplied literature source: \emph{Stability of Llarull's theorem in all dimensions}, \url{https://www.sciencedirect.com/science/article/abs/pii/S0001870824004961}.
\item Literature assessment dated exactly \texttt{2026-08-22}: the supplied curation labels the problem \texttt{partially\_solved} and reports that spinorial methods establish several positive-mass and scalar-curvature stability results, including spin and low-dimensional cases, while the requested all-dimensional Geroch and Goette--Semmelmann generalizations are not fully proved. This is supplied dated curation, not an independently verified literature assessment.
\item The registry labels its statement \texttt{corrected\_verified} and reports reconstruction and checking against the official Oberwolfach report. That verification claim is supplied provenance and has not been independently checked for this draft.
\end{itemize}
\end{document}
